\section{Relative ampleness on an arbitrary base and actual polarized base change}
\label{sec:polarized-base-change-v056}

\status{This chapter proves the full original
\texttt{relAmple\_of\_locally} for arbitrary schemes and the full
\texttt{PolBaseChangeSt}.  The construction uses the specified Cartesian
group-scheme pullback, keeps the original relative dimension, and equips
the actual pulled line bundle with its rigidification, symmetry,
relative ampleness and translation-kernel property.  The immediate
perfect-complex consumer also gives the full actual-section
finiteness statement and finiteness of every integer-power section
module.  There is no
Noetherianity or flatness hypothesis on the change of base.
The previously proved Serre and perfect-complex constructions are
reused at their existing scopes; the remaining multiplication,
geometric-rank, square-law and moduli producers are not asserted here.}

The main difficulty is to turn locally supplied ample presentations into
one presentation over each affine base open.  Their powers of the line
bundle can differ.  Taking a common positive multiple resolves that
problem, but merely extending powers of the original sections would
not yet prove the required generation of each affine coordinate ring.
We therefore extend all the mixed products
$s_b s_a^d$.  A finite ideal argument proves that their ratios generate
the nonvanishing open of every section of the common degree.
The argument works over arbitrary rings, including rings with
nilpotents.

The first part uses stalk-valued functions to express the construction
without constantly changing trivializations.  The second part recovers
actual compatible regular functions on every original line-bundle
chart.  Thus the auxiliary functions never replace the actual sections
required by the final statement.

\paragraph{Source and notation.}
Source commit:
\texttt{7cbf7a0c32fca1badf05931879a37c3f8a2b03b4}.
The principal files are
\texttt{Exists\allowbreak One\allowbreak Add\allowbreak Kill.lean},
\texttt{Rel\allowbreak Ample\allowbreak Of\allowbreak Locally\allowbreak On\allowbreak Base.lean}, and
\texttt{Pol\allowbreak Base\allowbreak Change\allowbreak Of\allowbreak Rel\allowbreak Ample\allowbreak Local.lean} in the original
abelian-scheme-law proofs.  The source map at the end identifies their
declarations and the additional chart and base-change arguments.

For a scheme $Z$ put
\[
 \operatorname{Fn}(Z)=\prod_{z\in Z}\mathcal O_{Z,z}.
\]
No regularity is required of an element of this product.  For an open
$W$, let $\operatorname{Reg}(W)$ be the subring consisting of the
functions which, on $W$, are the germs of one actual section of
$\mathcal O(W)$.  Values outside $W$ are immaterial.  Write
$\operatorname{inv}(F)(z)$ for the inverse of $F(z)$ if it is a unit
and for zero otherwise.  This is the source's total
$\operatorname{Ring.inverse}$, so identities involving it are used
only where the indicated element is a unit.

The assertion $\operatorname{IsLoc}(D,W,F)$ means
\[
 z\in D\quad\Longleftrightarrow\quad
 z\in W\ \hbox{and}\ F(z)\in\mathcal O_{Z,z}^{\times}.
                                                        \tag{PBC.1}
\]
For $S\subseteq\operatorname{Fn}(Z)$, write
$\operatorname{Gen}(D,S)$ if every $F\in\operatorname{Reg}(D)$
agrees on $D$ with an element of the subring $\langle S\rangle$.
This definition does not require the generators to be regular away
from $D$, and it does not identify functions which differ away from $D$.

Fix a cover $(U_i)_{i\in I}$ of $Z$ and functions
$G_i\in\operatorname{Fn}(Z)$ which are units at every point.  Assume
\[
 G_i\,\operatorname{inv}(G_j)\in
       \operatorname{Reg}(U_i\cap U_j).
                                                        \tag{PBC.2}
\]
For $n\ge0$ define
\[
 \operatorname{Sc}(W,n,F)
 \quad\Longleftrightarrow\quad
 G_i^nF\in\operatorname{Reg}(W\cap U_i)
       \quad\hbox{for every }i.
                                                        \tag{PBC.3}
\]
The cover and the functions $G_i$ remain fixed whenever this notation
is used.  They need not arise from a line bundle until PBC14.

\subsection{Denominator clearing and the finite ideal argument}

% BEGIN PBC1
\paragraph{PBC1: regular germs, unit loci and ratios.}
\label{claim:PBC1}
Equality of all germs on an open implies equality of its sections.
Regularity of a stalk-valued function is local.  A regular function
which is a unit in every stalk on $W$ has a regular inverse on $W$.
Consequently a function satisfying $\operatorname{Sc}(W,n,F)$
has an open unit locus in $W$.  If $F,\sigma$ both satisfy this
condition and $\operatorname{IsLoc}(D,W,\sigma)$, then
$F\operatorname{inv}(\sigma)\in\operatorname{Reg}(D)$.

\begin{proof}
The first assertion is the separatedness of the structure sheaf.
For the second, choose local representatives.  On an overlap their
germs agree with the prescribed function, so the representatives agree;
the sheaf gluing axiom gives one representative.  If $a\in\mathcal O(W)$
has unit germs, each point has a neighborhood on which $a$ has an
inverse: represent the inverse germ and shrink until its product with
$a$ is one.  These inverses agree on overlaps and glue.

On $W\cap U_i$ the function $G_i^nF$ is regular.  Since $G_i$ is a
unit, its ordinary principal open is exactly the locus where $F$ is
a unit.  The union of these opens is the asserted open.  On
$D\cap U_i$, both $G_i^nF$ and $G_i^n\sigma$ are regular and the
second is a unit.  Their quotient is
\[
 \frac{G_i^nF}{G_i^n\sigma}=F\operatorname{inv}(\sigma).
\]
These local representatives agree at every germ, hence glue.
The same definitions also give restriction of $\operatorname{Sc}$,
addition in one degree, multiplication with degrees added, and
$\operatorname{Sc}(W,nd,F^d)$ from
$\operatorname{Sc}(W,n,F)$.
\end{proof}
% END PBC1

% BEGIN PBC2
\paragraph{PBC2: affine fractions and one exponent on a quasi-compact open.}
\label{claim:PBC2}
Suppose $\operatorname{IsLoc}(D,W,w)$ and
$w\in\operatorname{Reg}(W)$.  If $W$ is affine and
$F\in\operatorname{Reg}(D)$, there are $a\ge0$ and
$H\in\operatorname{Reg}(W)$ such that
$w^aF=H$ on $D$.  If $W$ is only quasi-compact,
$F\in\operatorname{Reg}(W)$ and $F=0$ on $D$, there is $b\ge0$
such that $w^bF=0$ on all of $W$.

\begin{proof}
Represent $w$ by $r\in\mathcal O(W)$.  When $W=\operatorname{Spec}A$,
the open $D$ is $D(r)$ and
$\mathcal O(D)=A[r^{-1}]$.  Write the representative of $F$ as
$h/r^a$ and use the germs of $h$ on $W$, with arbitrary values
outside $W$, for $H$.

For the second assertion, choose a finite affine cover $V_\alpha$
of $W$.  The image of $F|_{V_\alpha}$ is zero after inverting
$w|_{V_\alpha}$, so some power of that element kills it in
$\mathcal O(V_\alpha)$.  The maximum of the finitely many exponents
works on every $V_\alpha$.  Its product with $F$ is therefore zero
on $W$.  Only the elementary equality criterion in a localization
is used; $w$ may be a zero divisor.
\end{proof}
% END PBC2

% BEGIN PBC3
\paragraph{PBC3: extending a prescribed-degree section across a principal base open.}
\label{claim:PBC3}
Let $Z$ be quasi-compact and quasi-separated, with the cover and
unit functions in (PBC.2).  Suppose $f\in\operatorname{Reg}(Z)$,
$\operatorname{IsLoc}(W,Z,f)$, and
$\operatorname{Sc}(W,n,F)$.  Then there are $k\ge0$ and
$\sigma\in\operatorname{Fn}(Z)$ such that
\[
 \operatorname{Sc}(Z,n,\sigma),\qquad
 \sigma=f^kF\quad\hbox{on }W.
                                                        \tag{PBC.4}
\]
The degree $n$ is unchanged.  This is the full original
$\operatorname{Sc.extend}$.

\begin{proof}
Choose finitely many affine opens $V_\alpha$ covering $Z$, each
contained in some $U_{i(\alpha)}$.  On
$W\cap V_\alpha$ the regular function
$G_{i(\alpha)}^nF$ has an affine fraction representation.  Thus
PBC2 gives $p_\alpha$ and $H_\alpha\in\operatorname{Reg}(V_\alpha)$
such that
\[
 H_\alpha=f^{p_\alpha}G_{i(\alpha)}^nF
       \quad\hbox{on }W\cap V_\alpha.
\]
Choose $p\ge p_\alpha$ for every $\alpha$ and put, as a function of
stalks,
\[
 t_\alpha=f^{p-p_\alpha}H_\alpha
                    \operatorname{inv}(G_{i(\alpha)})^n.
\]
Then $t_\alpha=f^pF$ on $W\cap V_\alpha$.  Although $t_\alpha$
itself need not be regular, for every $i$ the identity
\[
 G_i^nt_\alpha
  =(G_i/G_{i(\alpha)})^n f^{p-p_\alpha}H_\alpha
                                                        \tag{PBC.5}
\]
shows that it is regular in the $i$th gauge on $U_i\cap V_\alpha$.

The intersection $V_\alpha\cap V_\beta$ is quasi-compact by
quasi-separatedness.  There the regular difference
$G_{i(\alpha)}^n(t_\alpha-t_\beta)$ vanishes after restricting to
$W$.  PBC2 gives an exponent $q_{\alpha\beta}$ killing it on the
whole intersection.  Cancel the unit $G_{i(\alpha)}^n$.
Choose one $q$ at least all these finitely many exponents.  Then
$f^qt_\alpha=f^qt_\beta$ on every overlap.

For each $z$ choose $\alpha(z)$ containing $z$ and define
$\sigma(z)=f(z)^qt_{\alpha(z)}(z)$.  On $V_\alpha$ this agrees
with $f^qt_\alpha$ regardless of the pointwise choice.  Formula
(PBC.5), followed by PBC1, proves
$G_i^n\sigma\in\operatorname{Reg}(U_i)$ for every $i$.
On $W$ the definition gives $\sigma=f^{p+q}F$.  Thus $k=p+q$
works.  The first exponent clears fractions; the second makes the
different local extensions agree.  Neither is discarded.
\end{proof}
% END PBC3

% BEGIN PBC4
\paragraph{PBC4: the principal-open affineness criterion.}
\label{claim:PBC4}
Let $W$ be an open in a scheme and let $r_\lambda\in\mathcal O(W)$
generate the unit ideal.  If each $D_W(r_\lambda)$ is affine,
then $W$ is affine.

\begin{proof}
A finite subfamily already generates the unit ideal, because an
expression for $1$ is a finite sum.  Its principal opens cover $W$.
They are quasi-compact.  Their pairwise intersections are principal
opens in either affine member, hence affine and quasi-compact.
It follows that $W$ is quasi-compact and quasi-separated.

Put $A=\mathcal O(W)$.  For each $r=r_\lambda$, restriction induces
\[
 A[r^{-1}]\longrightarrow\mathcal O(D_W(r)).
                                                        \tag{PBC.6}
\]
It is injective by PBC2: a section restricting to zero is killed by a
power of $r$.  It is surjective by PBC3 in degree zero with the
trivial gauge $G=1$: a section on $D_W(r)$ becomes the restriction
of a global section after multiplication by a power of $r$.
Thus (PBC.6) is an isomorphism.

The canonical map $W\to\operatorname{Spec}A$ therefore restricts
to an isomorphism $D_W(r_\lambda)\simeq
D_{\operatorname{Spec}A}(r_\lambda)$ for the finite chosen family.
These target opens cover because the $r_\lambda$ generate $1$.
The restrictions agree on overlaps, being induced by the same
global ring map.  Their inverses glue to an inverse of the
canonical map.  Hence $W$ is affine.
\end{proof}
% END PBC4

% BEGIN PBC5
\paragraph{PBC5: generated functions, localization, and the exact polynomial map.}
\label{claim:PBC5}
The following three operations preserve the source's
$\operatorname{Gen}$ assertion.
First, generators may be replaced by elements agreeing on the
relevant open with polynomials in new generators.
Second, if $W$ is affine,
$\operatorname{Gen}(W,S)$ and
$\operatorname{IsLoc}(D,W,w)$ with $w\in\operatorname{Reg}(W)$,
then
\[
 \operatorname{Gen}(D,S\cup\{\operatorname{inv}(w)\}).
                                                        \tag{PBC.7}
\]
Third, let $\psi:R\to\mathcal O(D)$ be a ring map and
$r_j\in\mathcal O(D)$ an arbitrary family.  If
$\phi(a),\rho_j\in\operatorname{Fn}(Z)$ have on $D$ the germs
of $\psi(a),r_j$, respectively, then
\[
 \operatorname{Gen}\bigl(D,\phi(R)\cup\{\rho_j\}\bigr)
 \quad\Longleftrightarrow\quad
 R[X_j]\longrightarrow\mathcal O(D),\
 X_j\longmapsto r_j
 \ \hbox{is surjective}.
                                                        \tag{PBC.8}
\]

\begin{proof}
An element of a generated subring is obtained by finitely many
additions, negations and multiplications from the generators and
$0,1$.  Substituting the asserted new expressions proves the first
operation by induction on such an expression.

For (PBC.7), write $w^aF=H$ on $D$ using PBC2, then choose
$b\in\langle S\rangle$ agreeing with $H$ on $W$.
On $D$ the function $F$ agrees with
$\operatorname{inv}(w)^a b$.

For (PBC.8), an expression in $\phi(R)$ and the $\rho_j$ becomes a
polynomial over $R$ by the same induction; its evaluation has exactly
the indicated germs.  Applied to the germs of any section of
$\mathcal O(D)$, equality of germs proves surjectivity.
Conversely, a polynomial representative of the section representing
$F\in\operatorname{Reg}(D)$ supplies the element of the generated
subring required by $\operatorname{Gen}$.  There is no finiteness,
affineness, or Noetherianity hypothesis in (PBC.8).  The actual
coefficient map $\psi$ is retained.
\end{proof}
% END PBC5

% BEGIN PBC6
\paragraph{PBC6: local affine generation gives global affine generation.}
\label{claim:PBC6}
Let $B\subseteq\operatorname{Fn}(Z)$ be a subring with
$B\subseteq\operatorname{Reg}(W)$.  Let
\[
 K=\{b\in B:b(z)=0\text{ for every }z\in W\}.
                                                        \tag{PBC.9}
\]
Suppose $w_\lambda\in B$ satisfy $K+(w_\lambda)=B$.
For each $\lambda$ let
$\operatorname{IsLoc}(O_\lambda,W,w_\lambda)$, assume
$O_\lambda$ affine, and assume
$\operatorname{Gen}(O_\lambda,B\cup\{\operatorname{inv}(w_\lambda)\})$.
Then $W$ is affine and $\operatorname{Gen}(W,B)$.
The family of $\lambda$ may be infinite.

\begin{proof}
Assign to $b\in B$ the unique actual section with its germs on
$W$.  PBC1 shows that this defines a ring map
$\rho:B\to\mathcal O(W)$ with kernel $K$.
The images $\rho(w_\lambda)$ generate the unit ideal.  The
$O_\lambda$ are their principal opens, so PBC4 proves affineness
of $W$.

Fix $F\in\operatorname{Reg}(W)$ and define the ideal
\[
 J_F=\{b\in B:\text{ there is }b'\in B
                  \text{ with }bF=b'\text{ on }W\}.
                                                        \tag{PBC.10}
\]
It contains $K$.  On $O_\lambda$, local generation expresses
$F$ as a polynomial in elements of $B$ and $w_\lambda^{-1}$.
Clearing the finitely many denominators gives
$w_\lambda^aF=b'$ there, with $b'\in B$.
The difference is regular on $W$.  As $W$ is now affine and in
particular quasi-compact, PBC2 gives $b\ge0$ with
$w_\lambda^b(w_\lambda^aF-b')=0$ on $W$.
Hence $w_\lambda^{a+b}\in J_F$.

Thus $\sqrt{J_F}$ contains $K$ and all $w_\lambda$; it is the
unit ideal.  Consequently $J_F$ itself contains $1$.
The definition of $J_F$ says that $F$ agrees on $W$ with an
element of $B$.  This is the required global generation.
\end{proof}
% END PBC6

% BEGIN PBC7
\paragraph{PBC7: the finite one-plus-an-ideal annihilator.}
\label{claim:PBC7}
Let $B$ be any commutative ring, $u_t$ a finite family,
$I=(u_t)$, and $K$ any ideal.  Suppose that for each $t$ there is
$e_t\ge0$ such that
\[
 u_t^{e_t}\in K+I^{e_t+1}.
                                                        \tag{PBC.11}
\]
There are $c\in B$ and $d\ge0$ with
\[
 c-1\in I,\qquad c u_t^d\in K\quad\hbox{for every }t.
                                                        \tag{PBC.12}
\]
This is the full original \texttt{exists\_one\_add\_kill}.

\begin{proof}
If the family is empty, take $c=1,d=0$.  If some $e_t=0$,
(PBC.11) says $1\in K+I$.  Write $1=c+i$ with $c\in K,i\in I$
and take $d=0$.
Otherwise every $e_t\ge1$.  Put
$D=1+\sum_t(e_t-1)$.  Every monomial of degree $D$ in the
$u_t$ is divisible by some $u_t^{e_t}$; multiplying its
(PBC.11) relation by the remaining factors gives
\[
 I^D\subseteq K+I^{D+1}.
                                                        \tag{PBC.13}
\]
Consider the finitely generated $B$-module
$M=(I^D+K)/K$, generated by the degree-$D$ monomials.
Equation (PBC.13) says $M=IM$.
Choose generators $m_1,\ldots,m_r$ and write
$m_i=\sum_j a_{ij}m_j$ with every $a_{ij}\in I$.
For the matrix $A=(a_{ij})$, the adjugate identity applied to
$(1-A)(m_j)_j=0$ shows that
$c=\det(1-A)$ kills every generator, hence all of $M$.
Expanding the determinant gives $c-1\in I$.
Thus $cI^D\subseteq K$ and $d=D$ works.
If $M=0$, the same conclusion follows from the empty determinant
$c=1$.  This is a finite-matrix argument, so no Noetherian
assumption on $B$ or finite generation of $K$ is involved.
\end{proof}
% END PBC7

% BEGIN PBC8
\paragraph{PBC8: mixed ratios and the extra-section ideal estimates.}
\label{claim:PBC8}
Let $T$ be finite and suppose sections in the stalk notation
$s_j$ have degree $N$ on $W_0$.  Their unit loci $D_j$ cover
$W_0$, are affine, and satisfy
\[
 \operatorname{Gen}\bigl(D_j,S_0\cup\{s_m/s_j:m\in T\}\bigr),
 \qquad S_0\subseteq\operatorname{Reg}(W_0).
                                                        \tag{PBC.14}
\]
Let $\operatorname{Sc}(W_0,N(d+1),\tau)$ with $d\ge0$, and
let $W$ be its unit locus in $W_0$.  Put
\[
 u_{ij}=s_j s_i^d/\tau,\quad
 B=\langle S_0\cup\{u_{ij}:i,j\in T\}\rangle,\quad
 I=(u_{ij})\subseteq B,
                                                        \tag{PBC.15}
\]
where quotients mean the prescribed total inverses.
Then $B\subseteq\operatorname{Reg}(W)$, the opens
$D_j\cap W$ are affine, and, for $K$ as in (PBC.9),
each $u_{ij}$ satisfies the hypothesis (PBC.11).

\begin{proof}
The numerator in $u_{ij}$ has the same degree as $\tau$, so
PBC1 makes $u_{ij}$ regular on $W$.  On $D_j$ the function
\[
 v_j=\tau/s_j^{d+1}
\]
is regular.  Its unit locus is $D_j\cap W$, hence this
intersection is a principal open in the affine $D_j$.
On $D_j\cap W$ one has
\[
 v_ju_{jj}=1,\qquad
 u_{jj}(s_m/s_j)=u_{jm}.
                                                        \tag{PBC.16}
\]
Since the finite family $D_j\cap W$ covers $W$, this last open
is quasi-compact.

By (PBC.14), $v_j$ agrees on $D_j$ with an expression $b_j$
in $S_0$ and the ratios $s_m/s_j$.  Induction on this expression,
using the second identity in (PBC.16), gives an integer $e\ge0$
and $x\in I^e$ with
\[
 u_{jj}^e b_j=x\quad\hbox{on }D_j\cap W.
                                                        \tag{PBC.17}
\]
Here a generator from $S_0$ uses exponent zero; a ratio uses
exponent one and $u_{jm}\in I$; products add exponents.
For a sum with exponents $a,b$, multiply its two numerators
by $u_{jj}^b,u_{jj}^a$, respectively.  Both new numerators
belong to $I^{a+b}$.  This proves (PBC.17), including sums
with coefficients and negative signs.

The first identity in (PBC.16) gives
$u_{jj}^e-u_{jj}x=0$ on $D_j\cap W$.
This intersection is also the unit locus of $u_{jj}$ in $W$.
PBC2, applied to the regular difference on the quasi-compact
open $W$, gives $m\ge0$ such that
\[
 u_{jj}^{m+e}
       \equiv u_{jj}^{m+1}x \pmod K.
                                                        \tag{PBC.18}
\]
The right side belongs to $I^{m+e+1}$.  Thus (PBC.11)
holds for each diagonal $u_{jj}$.

For all $i,j$, the pointwise identity
\[
 u_{ij}^{d+1}=u_{jj}u_{ii}^{d}
                                                        \tag{PBC.19}
\]
holds even outside $W$: it is a polynomial identity in the
functions $s_i,s_j,\operatorname{inv}(\tau)$.
If the diagonal relation uses exponent $e$, raising
(PBC.19) to $e$ and multiplying that relation by $u_{ii}^{de}$
shows
$u_{ij}^{(d+1)e}\in K+I^{(d+1)e+1}$.
This proves the assertion for every mixed ratio.
\end{proof}
% END PBC8

% BEGIN PBC9
\paragraph{PBC9: every section of a higher common degree has an affine generated unit locus.}
\label{claim:PBC9}
Under the full hypotheses of PBC8,
\[
 W\text{ is affine},\qquad
 \operatorname{Gen}\bigl(W,S_0\cup\{s_j s_i^d/\tau:i,j\in T\}\bigr).
                                                        \tag{PBC.20}
\]
This proves the original \texttt{extra\_section} for every
$d\ge0$ and every such $\tau$.

\begin{proof}
Apply PBC7 to $I,K$ and the mixed ratios.  We obtain
$c-1\in I$ and one $q$ with $c u_{ij}^q\in K$ for every pair.
At a point of $W$, choose $j$ with the point in $D_j$.
Then $u_{jj}$ is a unit by (PBC.16).  Evaluating
$c u_{jj}^q=0$ at that point shows $c=0$ there.
Thus $c\in K$.

Let $J=K+(u_{jj}:j\in T)$.  Formula (PBC.19) implies
$I\subseteq\sqrt J$.  Since $c\in J$ and $c-1\in I$,
we have $1\in\sqrt J$, whence $J=B$.

We can now apply PBC6 with the diagonal functions $u_{jj}$.
Their unit loci in $W$ are the already affine opens
$D_j\cap W$.  To check its generation condition, first localize
(PBC.14) at $v_j$ using PBC5.  On this open,
\[
 v_j^{-1}=u_{jj},\qquad
 s_m/s_j=u_{jj}^{-1}u_{jm}.
                                                        \tag{PBC.21}
\]
The localized generators are therefore expressions in $B$ and
$u_{jj}^{-1}$.  PBC6 proves that $W$ is affine and generated
by $B$.  Finally, substituting the defining generators of $B$
gives precisely (PBC.20).
If $T$ is empty the cover forces $W_0$, and hence $W$, to be
empty; the same argument uses the zero coordinate ring and remains
valid.  No nonemptiness condition has been added.
\end{proof}
% END PBC9

\subsection{One positive power and one global ample pack}

A \emph{pack} on $W$ with base generators $S$ consists of a positive
integer $N$, a nonempty finite tuple $s_0,\ldots,s_m$ satisfying
$\operatorname{Sc}(W,N,s_j)$, and their unit loci $D_j$ in $W$,
such that these opens cover $W$, are affine, and satisfy
$\operatorname{Gen}(D_j,S\cup\{s_b/s_j\}_b)$.
This is exactly the source's $\operatorname{Pack}(U,G,W,S)$.

% BEGIN PBC10
\paragraph{PBC10: synchronizing all local degrees before extension.}
\label{claim:PBC10}
Let $Z$ be quasi-compact and quasi-separated, with (PBC.2).
Let $\phi:R\to\operatorname{Fn}(Z)$ be a ring map whose values
are regular on $Z$.  Let $(a_\ell)_{\ell\in\Lambda}$ be a finite
family generating the unit ideal of $R$, and let $W_\ell$ be
the unit locus of $\phi(a_\ell)$ in $Z$.
Suppose a pack is given on each $W_\ell$ with base generators
\[
 S_\ell=\phi(R)\cup\{\operatorname{inv}(\phi(a_\ell))\}.
                                                        \tag{PBC.22}
\]
There is a single $N\ge1$ and, for each
$t=(\ell,a,b)$ with $0\le a,b\le m_\ell$, a global
degree-$N$ function $\sigma_t$, together with $k_t\ge0$, such that
\[
 N=N_\ell(d_\ell+1),\qquad
 \sigma_{\ell,a,b}
    =\phi(a_\ell)^{k_{\ell,a,b}}\,
       s_{\ell,b}s_{\ell,a}^{d_\ell}
       \quad\hbox{on }W_\ell .
                                                        \tag{PBC.23}
\]
The choices of $N,\sigma_t,k_t$ are made before any global
degree-$N$ section later used as a denominator.

\begin{proof}
Choose all the finitely many local pack data and put
$N=\prod_{\ell\in\Lambda}N_\ell$.  This is positive, with
the empty product equal to $1$.  Each $N_\ell$ divides $N$
with positive quotient, which we write $d_\ell+1$.
Multiplication and powers in (PBC.3) show that
$s_{\ell,b}s_{\ell,a}^{d_\ell}$ has degree $N$ on
$W_\ell$.  Apply PBC3 to this function with
$f=\phi(a_\ell)$.  It gives exactly (PBC.23).
There are only finitely many triples, but no common bound on their
extension exponents is needed.

The $W_\ell$ cover $Z$.  Indeed, if at a point all
$\phi(a_\ell)(z)$ were nonunits, they would all belong to the
maximal ideal of the local ring $\mathcal O_{Z,z}$.
The image of the equation $1\in(a_\ell)$ would put $1$ in that
maximal ideal, a contradiction.
\end{proof}
% END PBC10

% BEGIN PBC11
\paragraph{PBC11: affine generation for every denominator of the common degree.}
\label{claim:PBC11}
Keep the data chosen in PBC10.  For every $\tau$ satisfying
$\operatorname{Sc}(Z,N,\tau)$, and every open $E$ with
$\operatorname{IsLoc}(E,Z,\tau)$, one has
\[
 E\text{ affine},\qquad
 \operatorname{Gen}\bigl(E,\phi(R)\cup\{\sigma_t/\tau\}_t\bigr).
                                                        \tag{PBC.24}
\]

\begin{proof}
Let $B=\langle\phi(R)\cup\{\sigma_t/\tau\}_t\rangle$.
The ratio assertion of PBC1 gives $B\subseteq\operatorname{Reg}(E)$.
There is a ring map $\psi:R\to B$ induced by $\phi$.
The elements $\psi(a_\ell)$ generate the unit ideal of $B$.

Apply PBC9 on $W_\ell$ to its original pack and the restriction
of $\tau$, using $N=N_\ell(d_\ell+1)$.
The set $S_\ell$ is regular there: the only additional generator
is the inverse of the unit $\phi(a_\ell)$.
It follows that $W_\ell\cap E$ is affine and generated by
$S_\ell$ and the functions
$s_{\ell,b}s_{\ell,a}^{d_\ell}/\tau$.
On this intersection (PBC.23) gives the exact identity
\[
 \frac{s_{\ell,b}s_{\ell,a}^{d_\ell}}{\tau}
   =\operatorname{inv}(\phi(a_\ell))^{k_{\ell,a,b}}
               \frac{\sigma_{\ell,a,b}}{\tau}.
                                                        \tag{PBC.25}
\]
Consequently these generators are expressions in $B$ and
$\operatorname{inv}(\psi(a_\ell))$.  The open
$W_\ell\cap E$ is the unit locus of $\psi(a_\ell)$ in $E$.
PBC6 now proves that $E$ is affine and generated by $B$.
PBC5 substitutes the generators of $B$ and gives (PBC.24).
This argument applies to every $\tau$ after the numerator family
has been fixed.
\end{proof}
% END PBC11

% BEGIN PBC12
\paragraph{PBC12: the full original global pack constructor.}
\label{claim:PBC12}
Under the original hypotheses of PBC10, there are $N\ge1$,
$M\ge0$, and $\sigma_0,\ldots,\sigma_M$ such that
$\operatorname{Sc}(Z,N,\sigma_t)$ for every $t$, at every point
some $\sigma_t$ is a unit, and for every $t$ and every open $D$
\[
 \operatorname{IsLoc}(D,Z,\sigma_t)
 \ \Longrightarrow\
 D\text{ affine and }
 \operatorname{Gen}\bigl(D,\phi(R)\cup\{\sigma_m/\sigma_t\}_m\bigr).
                                                        \tag{PBC.26}
\]
This is the full original \texttt{glue\_pack}, including its
universal final open $D$.

\begin{proof}
Adjoin one zero function to the finite family in PBC10 and index
the resulting nonempty family by
$\{0,\ldots,M\}$, where $M$ is the number of triples
$(\ell,a,b)$.  Zero satisfies $\operatorname{Sc}$ in every degree.
At a point $z$, choose $\ell$ with $z\in W_\ell$ and then a local
section $s_{\ell,j}$ which is a unit at $z$.  The diagonal case
of (PBC.23) gives
\[
 \sigma_{\ell,j,j}(z)
   =\phi(a_\ell)(z)^{k_{\ell,j,j}}
                     s_{\ell,j}(z)^{d_\ell+1},
\]
which is a unit.  Thus the unit loci cover.

For the last assertion apply PBC11 with $\tau=\sigma_t$.
All its numerator functions occur in the enlarged family,
so PBC5 gives the displayed generation statement.
If the initial finite index set is empty, the unit-ideal
hypothesis and the local-ring argument of PBC10 force $Z$ to be
empty; the appended zero still gives the required nonempty tuple.
This accounts for the source's $\operatorname{Fin}(M+1)$
without imposing an extra nonemptiness hypothesis.
\end{proof}
% END PBC12

\subsection{Returning the pack to the actual line-bundle charts}

% BEGIN PBC13
\paragraph{PBC13: transport through an open immersion.}
\label{claim:PBC13}
For a scheme map $h:Y\to Z$, define
\[
 (h^\flat F)(y)=h_y^\#(F(h(y))).
                                                        \tag{PBC.27}
\]
This is a ring map, composes as pullback, preserves and reflects
the property of being a unit at the corresponding stalk, and
commutes with the total inverse.  It preserves regularity by
pulling actual sections.  If $h$ is an open immersion, every
function in $\operatorname{Fn}(Y)$ has a lift in
$\operatorname{Fn}(Z)$, and regularity and generation on an open
of $Y$ transport to its actual image in $Z$.
In particular, a pack for the pulled cover, pulled gauges and
pulled generators on $Y$ gives a pack on $h(Y)$ for the original
cover, gauges and generators.

\begin{proof}
The stalk maps of a scheme morphism are local homomorphisms.
A local homomorphism preserves units and takes nonunits to nonunits.
Therefore it commutes with inverse in both cases of the definition.
The ring and composition assertions follow from the stalk maps,
and a pulled section has exactly the germs in (PBC.27).

For an open immersion, $h$ identifies $Y$ with its image and
every stalk map is an isomorphism.  Given $F$ on $Y$, use the
inverse stalk maps to define its lift on $h(Y)$ and define the
lift to be zero elsewhere.  This is a lift as a function of
stalks; it is not claimed to be regular across the boundary.
Regularity on an open $V\subseteq Y$ transports because the
open-immersion section map identifies $\mathcal O(h(V))$ with
$\mathcal O(V)$.  To transport $\operatorname{Gen}$, pull a
regular function back, use its generating expression there, and
use the injective stalk maps to obtain the same expression on
$h(V)$.

Now lift each of the finitely many functions in the given pack.
On $h(Y)\cap U_i$, the regularity of its $i$th gauge follows
from regularity on $Y\cap h^{-1}U_i$ and the preceding section
identification.  Unit loci are images of the old unit loci,
so they cover and are affine.  Pullback commutes with their
ratios, and the generation assertion just proved transports
the old ratio generators to the lifted ones.  These are all
fields of the asserted pack.
\end{proof}
% END PBC13

% BEGIN PBC14
\paragraph{PBC14: chosen-chart values and the literal coordinate-ratio identity.}
\label{claim:PBC14}
Let $P$ be an original line-bundle chart datum on $Y$.
Its transition convention for a section is
$s_k=g_{kk'}s_{k'}$ on $P.U_k\cap P.U_{k'}$.
Choose a chart $\gamma(y)$ containing each $y$, and write
\[
 \operatorname{sv}_\gamma(s)(y)
       =\operatorname{germ}_y(s_{\gamma(y)}).
                                                        \tag{PBC.28}
\]
For every actual section $s$, its original nonzero locus is the
unit locus of $\operatorname{sv}_\gamma(s)$.  For every actual
section $t$ and $r\in\mathcal O(D(s))$, the identities
\[
 t_k|_{D(s)\cap P.U_k}
   =r|_{D(s)\cap P.U_k}\,s_k|_{D(s)\cap P.U_k}
       \quad\hbox{for every }k
                                                        \tag{PBC.29}
\]
are equivalent to
$\operatorname{germ}_y(r)=
\operatorname{sv}_\gamma(t)(y)/
\operatorname{sv}_\gamma(s)(y)$ for every $y\in D(s)$.

\begin{proof}
At any point in two charts the two coordinates of $s$ differ
by the unit germ of $g_{kk'}$.  One is a unit if and only if
the other is.  Taking the union of their ordinary principal
opens is therefore the unit locus in (PBC.28), exactly the
source's $\operatorname{nonzeroLocus}$.

Equation (PBC.29) in the chosen chart gives the asserted ratio
after division by the unit coordinate of $s$.
Conversely, multiply the ratio equality by that coordinate,
then multiply both sides by the same transition unit to obtain
the equality in any other chart $k$ containing the point.
Equality of all germs proves (PBC.29).
Only the given section compatibility and the chosen chart cover
are needed for these two assertions.

More generally, let $U_i\subseteq P.U_{\kappa(i)}$ and choose
$c(y)$ with $y\in U_{c(y)}$.  Suppose on $U_i$ that
\[
 \operatorname{germ}_y g_{\kappa(i),\kappa(c(y))}
      =G_i(y)^N.
                                                        \tag{PBC.30}
\]
Then every actual section $s$ gives
$\operatorname{Sc}(Y,N,\operatorname{sv}_{\kappa c}(s))$:
its $i$th gauged function is exactly the germ of
$s_{\kappa(i)}|_{U_i}$.  This is the forward
$\operatorname{sc\_sv}$ dictionary.
\end{proof}
% END PBC14

% BEGIN PBC15
\paragraph{PBC15: reconstructing an actual section from every gauged function.}
\label{claim:PBC15}
Assume $P.\operatorname{IsLB}$ and the precise data
$U_i,\kappa,c,G_i,N$ of (PBC.30).  For this assertion the
$G_i$ can be arbitrary stalk functions: no additional
global-unit or regular-ratio hypothesis is needed.
Every $\sigma$ with
$G_i^N\sigma\in\operatorname{Reg}(U_i)$ for all $i$ is
$\operatorname{sv}_{\kappa c}(s)$ for an actual section
$s$ on every original chart of $P$.
This is the full original \texttt{exists\_sv}.

\begin{proof}
For each original chart $k$, prescribe on $P.U_k$ the stalk function
\[
 F_k(y)=
   \operatorname{germ}_y g_{k,\kappa(c(y))}\,\sigma(y).
                                                        \tag{PBC.31}
\]
Near a fixed $z\in P.U_k$, use the fixed index $i=c(z)$.
On $P.U_k\cap U_i$, the cocycle identity and (PBC.30) give
\[
 F_k=
   \operatorname{germ}(g_{k,\kappa(i)})\,G_i^N\sigma.
                                                        \tag{PBC.32}
\]
The first factor is the germ of an actual transition unit on
this open, and the second is regular there by hypothesis.
Thus $F_k$ is locally regular at every point of $P.U_k$.
PBC1 gives one actual $s_k\in\mathcal O(P.U_k)$ with these germs.

On $P.U_k\cap P.U_{k'}$, the cocycle equation
$g_{k,\kappa(c(y))}=g_{kk'}g_{k',\kappa(c(y))}$
shows that $s_k=g_{kk'}s_{k'}$ at every germ, hence as
actual sections.  The family $(s_k)$ is therefore an
$\operatorname{LB.Sec}(P)$ in the original coordinates.
At $k=\kappa(c(y))$ the transition is $1$, so (PBC.31)
gives $\operatorname{sv}_{\kappa c}(s)=\sigma$.
The argument glues sections on their actual algebraic opens;
it does not assign a geometric regularity to the pointwise
choice $c$.
\end{proof}
% END PBC15

% BEGIN PBC16
\paragraph{PBC16: the exact base constants on a principal neighborhood.}
\label{claim:PBC16}
Let $B$ be affine, $a\in\mathcal O(B)$, and let the open
immersion $b:V\to B$ have image $D_B(a)$.  Suppose
$q:Y\to V$, $h:Y\to Z$, and $p:Z\to B$ satisfy
$p\circ h=b\circ q$.  Then every pulled base function
$q^*c$, $c\in\mathcal O(V)$, is, as a function of all
stalks of $Y$, an expression in
\[
 h^\flat(p^*\mathcal O(B))
       \quad\hbox{and}\quad
 h^\flat\bigl(\operatorname{inv}(p^*a)\bigr).
                                                        \tag{PBC.33}
\]

\begin{proof}
The actual open immersion identifies
$\mathcal O(V)$ with $\mathcal O(B)[a^{-1}]$.
Choose $r\in\mathcal O(B)$ and $n\ge0$ representing $c=r/a^n$.
Pullback gives
$q^*c=(p\circ h)^*r\,((p\circ h)^*a)^{-n}$.
The denominator is a unit everywhere on $Y$, since $q$ factors
through $D_B(a)$.  PBC13 therefore identifies its inverse with
the total inverse appearing in (PBC.33).
This proves literal membership in the generated subring of
$\operatorname{Fn}(Y)$, not just agreement on a smaller open.
\end{proof}
% END PBC16

% BEGIN PBC17
\paragraph{PBC17: a local relative-ample presentation gives the required principal pack.}
\label{claim:PBC17}
Let $A\to S$, an original chart datum $L$, an affine open
$U\subseteq S$, and an open $U_0\subseteq S$ be given.
Suppose the actual restriction of $L$ to $A_{U_0}$ is
relatively ample in the original sense.
Let $a\in\mathcal O(U)$ satisfy
$D_U(a)\subseteq U\cap U_0$.
Put $Z=A_U$, $p:Z\to U$, $j:Z\to A$, and choose $c(z)$ with
$j(z)\in L.U_{c(z)}$.
Set
\[
 U_i=j^{-1}(L.U_i),\qquad
 G_i(z)=\operatorname{germ}_z(j^*g_{i,c(z)})
                 \quad(z\in U_i),
                                                        \tag{PBC.34}
\]
using the source's value $1$ for this unit expression when the
indicated overlap condition fails.
There is a pack on $p^{-1}D_U(a)$ in these gauges with base
generators $p^*\mathcal O(U)\cup\{\operatorname{inv}(p^*a)\}$.
No additional global ampleness hypothesis is used.

\begin{proof}
View $D_U(a)$ also as the open $V$ of $U_0$ given by its image
in $S$.  This is an affine open, and its map $b:V\to U$
is the specified identification with $D_U(a)$.
Let $Y=(A_{U_0})_V$ with projection $q:Y\to V$.
There is a unique open immersion $h:Y\to A_U$ commuting with
the maps to $A$; its image is exactly $p^{-1}D_U(a)$ and
$p\circ h=b\circ q$.  These statements follow directly by
restricting the common point and structure-sheaf maps to the
same open of $A$.  No new model of $A$ is chosen.

Apply the given local relative ampleness on the affine open
$V\subseteq U_0$.  It supplies $N\ge1$, sections
$s_0,\ldots,s_m$ of the original repeated-factor power on
$Y$, affine nonzero loci covering $Y$, and ratio functions
whose polynomial evaluation over $\mathcal O(V)$ is surjective.
For this power use constant chart tuples: the same original
$L$-chart is used in each of the $N$ factors.
Its transition on two such tuples is the $N$th power of
the original transition, because its definition is the product
of those $N$ identical factors.  Thus (PBC.30) holds for the
pulled gauges $h^\flat G_i$ and the choice $c(h(y))$.

PBC14 sends each actual $s_j$ to its chosen-chart value,
with exactly its old nonzero locus.  The coordinate-ratio
identities and PBC5 turn the supplied polynomial surjections
into $\operatorname{Gen}$ with the base generators
$q^*\mathcal O(V)$.
PBC16 expresses every one of these base generators in the
pulled family
$h^\flat(p^*\mathcal O(U)\cup\{\operatorname{inv}(p^*a)\})$.
Replacing generators by PBC5 gives a pack with this latter
base family.  Finally PBC13 transports it through $h$.

For a point $z\in Z$, the local stalk map of $p$ reflects
units.  Hence $p^*a$ is a unit at $z$ precisely when
$p(z)\in D_U(a)$.  The pack open just constructed is
therefore the literal unit locus required in the source.
The forward dictionary used here requires only the supplied
sections and chart choice; the assertion has the original
\texttt{pack\_of\_relAmple} scope.
\end{proof}
% END PBC17

% BEGIN PBC18
\paragraph{PBC18: full relative ampleness from neighborhoods of every base point.}
\label{claim:PBC18}
Let $S$ be an arbitrary scheme, $A\to S$ a quasi-compact,
quasi-separated morphism, and $L$ an actual line bundle on
$A$.  Suppose that for every $s\in S$ there is an open
$U_0$ containing $s$ such that the actual restriction
$L|_{A_{U_0}}$ is relatively ample.  Then
$\operatorname{RelAmple}(A,L)$ holds with the exact source
definition: for every affine open $U\subseteq S$ there are
$N\ge1$ and $m+1$ actual sections of
$\operatorname{powOver}(A,L,U,N)$ whose nonzero loci cover,
are affine, and have the stated coordinate-ratio polynomial
surjections over $\mathcal O(U)$.

\begin{proof}
Fix the affine open $U$ and put $Z=A_U$.
It is quasi-compact because $A\to S$ is quasi-compact.
It is quasi-separated because the restricted morphism is
quasi-separated and the affine base $U$ is separated.
Choose for every $z\in Z$ an original chart $c(z)$ of $L$
containing its image, and use (PBC.34).
These $G_i$ are units everywhere by their definition.
On $U_i\cap U_j$ the cocycle identity gives
\[
 G_i/G_j=j^*g_{ij}.
                                                        \tag{PBC.35}
\]
Thus the hypotheses (PBC.2) hold.

Call $a\in\mathcal O(U)$ good if $D_U(a)$ lies in a
neighborhood on which the assumed relative ampleness holds.
The principal-open basis on the affine scheme $U$ shows
that the $D_U(a)$ for good $a$ cover $U$:
at each point take its given neighborhood and then a
principal neighborhood contained in it.
Consequently the good elements generate the unit ideal.
Indeed, a proper ideal containing all of them would lie
in a prime ideal, whose point would be outside that cover.
Choose a finite expression for $1$ in this ideal.
It yields finitely many good $a_\ell$ which still generate $1$.
No finite-generation theorem for arbitrary ideals is being used.

For $\phi:\mathcal O(U)\to\operatorname{Fn}(Z)$ use the
actual pullback of base functions.  PBC17 gives the pack
required over each unit locus of $\phi(a_\ell)$.
Apply PBC12 to obtain one $N\ge1$ and global functions
$\sigma_0,\ldots,\sigma_m$ of degree $N$.

Now set $P=\operatorname{powOver}(A,L,U,N)$.
For $i$ let $\kappa(i)$ be the constant tuple of $N$
copies of $i$.  The original $L$-cocycle gives $P.\operatorname{IsLB}$;
each of its transition units is the product of the corresponding
pulled transition units.  On the constant tuples its transition
is $G_i^N$.  Hence PBC15 gives an actual section $s_t$ of
$P$ with chosen-chart value exactly $\sigma_t$.
PBC14 identifies its original nonzero locus with the unit
locus of $\sigma_t$.  PBC12 therefore gives the covering
and the affineness of every such open.

Fix $i$.  PBC1 gives a regular function
$r_{ij}$ on $D(s_i)$ whose germs are $\sigma_j/\sigma_i$.
By PBC14, on every original $P$-chart it satisfies
$s_j=r_{ij}s_i$.  The generation assertion (PBC.26),
combined with PBC5 for the actual coefficient map
$\mathcal O(U)\to\mathcal O(D(s_i))$, gives the surjection
\[
 \mathcal O(U)[X_0,\ldots,X_m]\longrightarrow
       \mathcal O(D(s_i)),\qquad X_j\longmapsto r_{ij}.
                                                        \tag{PBC.36}
\]
These are all fields in the original definition of
$\operatorname{RelAmple}$.  Since $U$ was arbitrary, the
assertion is the full \texttt{relAmple\_of\_locally}.
\end{proof}
% END PBC18

\subsection{Affine base change with the prescribed bundle and ratios}

% BEGIN PBC19
\paragraph{PBC19: generated bundle comparisons, powers, and actual pulled sections.}
\label{claim:PBC19}
Let $F:X'\to X$, let $L,L'$ be actual line bundles, and let
$w$ be a generated comparison $L'\simeq F^*L$ in the source
sense.  It gives a pulled section of $L'$ for every actual
section of $L$.  This section satisfies the comparison at
every point of every test scheme, its nonzero locus is the
inverse image of the original one, and ratios of two such
sections are the pulled ratios.
The same comparisons exist for every finite
$\operatorname{LinComb}$ with arbitrary integer exponents,
and for the literal nested pullback and composed pullback.

\begin{proof}
On $F^{-1}(L.U_j)\cap L'.U_k$, write $w_{jk}$ for the
specified unit in the division bundle.  Its compatibility is
\[
 w_{jk}=F^*g_{jj'}(g'_{kk'})^{-1}w_{j'k'}
                                                        \tag{PBC.37}
\]
on each common overlap.  For a section $s$ of $L$, define
locally on this open
$s'_k=w_{jk}^{-1}F^*s_j$.
When $j$ changes and $k$ stays fixed, (PBC.37) and
$s_j=g_{jj'}s_{j'}$ show that these functions agree.
They glue to an actual function on $L'.U_k$.
When $k$ changes, the same identity gives
$s'_k=g'_{kk'}s'_{k'}$.
Thus they form an actual section of $L'$.
Every equality is an equality of regular functions, so after
pullback along any $r:Z\to X'$ it remains the exact germ
identity
\[
 \operatorname{germ}(F\circ r)^*s_j
    =\operatorname{germ}(r^*w_{jk})\,
                      \operatorname{germ}(r^*s'_k).
                                                        \tag{PBC.38}
\]
This is the source's all-test \texttt{sec\_pull\_down}
conclusion.  A unit multiplier preserves the property of a
germ being a unit, and a local stalk homomorphism reflects
units.  The asserted nonzero-locus equality follows.
For two sections the same unit multiplies both coordinates,
so it cancels from their ratio.

For maps $\alpha_\ell:Y\to X'$ and integers $\epsilon_\ell$,
on a common refinement multiply the pulled units as
\[
 \prod_\ell \alpha_\ell^*(w_{j_\ell k_\ell})^{\epsilon_\ell}.
                                                        \tag{PBC.39}
\]
Negative exponents use inverses of units.
Multiplying (PBC.37) to these powers gives exactly the ratio
of the two $\operatorname{LinComb}$ transition functions.
Hence (PBC.39) is a generated comparison on the original
refined charts.  It carries a generating tuple to the
product of that tuple with (PBC.39), proving the
generator-transport statement as well.
For completeness, fix a chart of the target
$\operatorname{LinComb}$ and intersect it with all charts
of the source.  They cover that target chart, since each
line bundle covers and the product index is finite.
On these intersections the prescribed product units agree,
by (PBC.37) and the old generator's compatibility.
Glue the units and their inverses to obtain a unit on the
whole target chart.  These units satisfy its transition
identity by the same equation.  After every test map $r$,
their germs equal the product (PBC.39) times the germs of
the old generator.  This is also the full all-test formula
used by the source's \texttt{exists\_up}.

For nested versus composed pullbacks, restriction and pullback
of the same original coordinate function agree by the
composition rule for the structure sheaf.  On matching
refined charts use the unit $1$ and on nonmatching charts
use the corresponding original transition.  The cocycle
identity proves compatibility and gives the comparison
between the literal chart records.  In particular the
identity comparison on a bundle $M$ is represented, on
the pair chart $(i,j)$ of its division bundle, by $g^M_{ij}$:
the identity
$g^M_{ij}=(g^M_{ii'}(g^M_{jj'})^{-1})g^M_{i'j'}$
is exactly its required compatibility.
Pulling any generating tuple along a chart refinement is
proved by this same equality and needs no new existence theorem.
\end{proof}
% END PBC19

% BEGIN PBC20
\paragraph{PBC20: relative ampleness when the base map lands in one affine open.}
\label{claim:PBC20}
Suppose the square
\[
 \begin{array}{ccc}
 A'&\xrightarrow{\phi}&A\\
 p'\downarrow&&\downarrow p\\
 S'&\xrightarrow{f}&S
 \end{array}
                                                        \tag{PBC.40}
\]
is Cartesian.  Let $L,L'$ be actual line bundles, with a
specified generated comparison $L'\simeq\phi^*L$.
Assume $\operatorname{RelAmple}(A,L)$, and let $U\subseteq S$
be an affine open containing $f(S')$.
Then $\operatorname{RelAmple}(A',L')$.
Here $S'$ and $f$ are arbitrary; $f$ need not be flat.
This is the full original \texttt{relAmple\_pull}.

\begin{proof}
Fix any affine open $U'\subseteq S'$.  Restrict (PBC.40)
to obtain the actual Cartesian square
\[
 \begin{array}{ccc}
 A'_{U'}&\xrightarrow{\psi}&A_U\\
 \downarrow&&\downarrow\\
 U'&\xrightarrow{f|_{U'}}&U .
 \end{array}
                                                        \tag{PBC.41}
\]
The specified inclusion $f(U')\subseteq U$ defines this
restricted map.  The universal property of (PBC.40) and
the two open restrictions proves (PBC.41) is Cartesian.

Apply relative ampleness on $U$.  It provides $N\ge1$
and actual sections $s_0,\ldots,s_m$ of
$\operatorname{powOver}(A,L,U,N)$ with affine nonzero
loci $D_i$, ratio functions $r_{ij}$, and surjections
\[
 R[X_0,\ldots,X_m]\twoheadrightarrow Q_i,
 \quad R=\mathcal O(U),\quad Q_i=\mathcal O(D_i).
                                                        \tag{PBC.42}
\]
PBC19 compares their power bundle with the literal
$\operatorname{powOver}(A',L',U',N)$.
Pull each section through this comparison, obtaining $s'_i$.
Its nonzero locus is
$D'_i=\psi^{-1}D_i$, so the new loci cover.

Since $U,U'$ are affine, $U'\to U$ is affine, and its
base change $\psi$ is affine.  Hence each $D'_i$ is
affine.  More explicitly, put $R'=\mathcal O(U')$.
The actual affine square on $D_i$ identifies
\[
 \mathcal O(D'_i)\simeq R'\otimes_R Q_i.
                                                        \tag{PBC.43}
\]
This is the defining affine fiber-product construction:
maps out of the tensor product are pairs of ring maps
agreeing on $R$, which is precisely the square's universal
property.  Thus (PBC.43) is valid without any flatness
assumption.

Set $r'_{ij}=\psi^*r_{ij}$.  PBC19 gives on every original
power-bundle chart the equality $s'_j=r'_{ij}s'_i$.
Every pure tensor in (PBC.43) has its $Q_i$ factor represented
by a polynomial from (PBC.42).  It is therefore a polynomial
over $R'$ in the images of the $r_{ij}$.  Sums of pure
tensors give all of the tensor product.  Hence
\[
 R'[X_0,\ldots,X_m]\twoheadrightarrow\mathcal O(D'_i),
       \qquad X_j\longmapsto r'_{ij}.
                                                        \tag{PBC.44}
\]
The coefficient map and the ratio maps are the actual maps
of (PBC.41).  We have produced every field of relative
ampleness on $U'$, with the same $N$ and the pulled tuple.
As $U'$ was arbitrary, the result follows.
\end{proof}
% END PBC20

% BEGIN PBC21
\paragraph{PBC21: the literal restricted pullback bundle is locally relatively ample.}
\label{claim:PBC21}
In (PBC.40), suppose $L$ is an actual relatively ample
line bundle on $A$.  For every affine open $V\subseteq S$,
the bundle
\[
 (\phi^*L)|_{A'_{f^{-1}V}}
                                                        \tag{PBC.45}
\]
with the literal nested pullback charts is relatively ample
over $f^{-1}V$.

\begin{proof}
Let $i:A'_{f^{-1}V}\to A'$ and
$b:f^{-1}V\to S'$ be the actual open immersions.
Pasting the restriction square with (PBC.40) gives the
Cartesian square with upper map $\phi\circ i$ and lower
map $f\circ b$.  Its base image is contained in $V$.
PBC19 supplies a generated comparison
$i^*(\phi^*L)\simeq(\phi\circ i)^*L$ between the literal
two chart records.  Apply PBC20 with this comparison and
the affine open $V$.
The result is exactly (PBC.45); the possibly nonaffine
open $f^{-1}V$ has not been replaced by a smaller base.
\end{proof}
% END PBC21

\subsection{The actual group pullback and all the polarization fields}

For the remaining assertions, $A/S$ and $A'/S'$ are group
objects, and a \emph{group pullback} means (PBC.40) together
with the original all-test multiplication assertion.
Explicitly, for every $T\to S'$ and maps $x,y:T\to A'$,
composition with $\phi$ identifies their product with the
product of their two images in $A(T\to S'\to S)$.
The notation $\operatorname{toBase}$ below means exactly
this composition, with the induced structure map on $T$.

% BEGIN PBC22
\paragraph{PBC22: the pullback group, relative dimension, and the exact test maps.}
\label{claim:PBC22}
For arbitrary $f:S'\to S$, the specified scheme
$A'=A\times_S S'$ has its canonical pulled group structure
and makes (PBC.40) a group pullback.  If $A/S$ is an
abelian scheme, so is $A'/S'$.  If its original smooth
relative dimension is $g$, the pullback has the same
smooth relative dimension $g$.
For any group pullback, not only this chosen model,
$\operatorname{toBase}$ is injective on every test scheme,
preserves identities and inverses, and
\[
 \phi\circ e_{A'}=e_A\circ f,\qquad
 \phi\circ[-1]_{A'}=[-1]_A\circ\phi .
                                                        \tag{PBC.46}
\]

\begin{proof}
For a test $T\to S'$, the universal property identifies
$A'(T)$ with the maps $a:T\to A$ whose composition with
$p$ is the prescribed map $T\to S'\to S$.
This identification is natural in $T$.  Define multiplication,
identity and inverse by those of $A(T)$; their structure
maps to $S'$ remain the prescribed one.
Equivalently, lift the original multiplication and inverse
through the Cartesian square.  Uniqueness of these lifts
proves the group axioms and the original all-test product
identity.  The given categorical pullback group has exactly
these lifts, so the construction is the original
$\operatorname{Over.pullback}$ group.

For any group pullback, two maps $T\to A'$ with the same
composite to $A$ also have the same composite to $S'$
because they are over $S'$.  Cartesian uniqueness gives
their equality.  Thus $\operatorname{toBase}$ is injective.
It is multiplicative by the supplied all-test field.
A multiplicative map of groups preserves the identity
and inverse: cancel its image of $1$ in
$h(1)=h(1)^2$, then apply $h(xx^{-1})=1$.
Apply this to the identity test and to the universal point
of $A'$ to obtain (PBC.46).

Commutativity of the original group law is preserved by
the same lifts.  Smoothness and its relative dimension
are preserved by base change: on a standard smooth
presentation, base change replaces the coefficients by
their images, leaves the number of variables and relations
unchanged, and preserves the invertible Jacobian minor.
Such presentations cover the pulled scheme.
Properness is also preserved: finite type and separatedness
are stable under base change, and universal closedness
holds because every further base change is already a base
change of the original proper morphism.
Finally, a geometric fiber of $A'\to S'$ over an
algebraically closed field is the same fiber of $A\to S$
after composition of the field point with $f$.
It is connected by the original geometric-connectedness
field.  These are precisely the four fields
commutative, smooth, proper, and geometrically connected
in the original $\operatorname{IsAbelianScheme}$ record.
\end{proof}
% END PBC22

% BEGIN PBC23
\paragraph{PBC23: pulling the prescribed rigidification and local symmetry.}
\label{claim:PBC23}
Let (PBC.40) be any group pullback.
If $L$ is rigidified along $e_A$, then the literal $\phi^*L$
is rigidified along $e_{A'}$.
If $L$ satisfies the original locally-on-$S$ symmetry
condition, then $\phi^*L$ satisfies the original
locally-on-$S'$ symmetry condition.
The actual generating tuples are pulled from the given ones.

\begin{proof}
The given rigidification is a compatible unit tuple
trivializing $e_A^*L$.  Pull it by $f$ and use the first
identity in (PBC.46) to compare
$f^*e_A^*L$ with $e_{A'}^*\phi^*L$.
On each common chart the tuple is the pullback of the old
unit.  Its compatibility follows by pulling the old
compatibility equation.  The chart comparison of PBC19
puts it in the literal target chart record.

For symmetry, fix $s'\in S'$.  At $f(s')$ the original
condition supplies an open $U\subseteq S$ and a
generating tuple of
$[-1]_A^*L\otimes L^{-1}$ on $A_U$.
Use the open $f^{-1}U$, which contains $s'$.
The Cartesian square restricts to a map
$A'_{f^{-1}U}\to A_U$.
The second identity in (PBC.46) says this map commutes
with the inversion maps.  Pull the given tuple and
apply the finite two-factor comparison (PBC.39) with
exponents $(1,-1)$.  It is a generator of precisely the
source's $\operatorname{LinComb}$ on this restricted
space.  This is the required local symmetry at $s'$.
No global symmetry generator is inferred from the local ones.
\end{proof}
% END PBC23

% BEGIN PBC24
\paragraph{PBC24: the full affine-test translation-kernel condition survives base change.}
\label{claim:PBC24}
Let (PBC.40) be a group pullback, let $L,L'$ be actual
line bundles, and fix a generated comparison
$L'\simeq\phi^*L$.
If $\operatorname{KTriv}(A,L)$ holds, then
$\operatorname{KTriv}(A',L')$ holds, for every commutative
test ring in its original definition.

\begin{proof}
Take any commutative ring $R$, put $T=\operatorname{Spec}R$,
and fix $t:T\to S'$ and $x:T\to A'$ over $S'$.
Write $x_0=\phi\circ x$, regarded over $f\circ t$.
The two spaces on which the translation differences are defined are
\[
 Y'=T\times_{S'}A',\qquad Y=T\times_S A.
                                                        \tag{PBC.47}
\]
There is a canonical map
$\alpha:Y'\to Y$ with first component $\operatorname{pr}_T$
and second component $\phi\circ\operatorname{pr}_{A'}$.
It is an isomorphism.  To construct its inverse, on $Y$ lift
the two compatible maps
$\operatorname{pr}_A:Y\to A$ and
$t\circ\operatorname{pr}_T:Y\to S'$ through (PBC.40), and
then pair the resulting map to $A'$ with $\operatorname{pr}_T$.
Both composites are identities because their two projections
are identities.  This proves the assertion on the actual
fiber-product models.

Let $\tau_x:Y'\to A'$ be addition of $x$ to the second
component and let $\tau_{x_0}:Y\to A$ be the analogous map.
The all-test multiplication field, applied to the universal
points on $Y'$, gives
\[
 \phi\circ\tau_x=\tau_{x_0}\circ\alpha,\qquad
 \phi\circ\operatorname{pr}_{A'}
                         =\operatorname{pr}_A\circ\alpha .
                                                        \tag{PBC.48}
\]
By PBC19 the specified comparison for $L',L$ therefore
gives a generated comparison
\[
 \tau_x^*L'\otimes\operatorname{pr}_{A'}^*(L')^{-1}
 \simeq
 \alpha^*\bigl(
    \tau_{x_0}^*L\otimes\operatorname{pr}_A^*L^{-1}\bigr).
                                                        \tag{PBC.49}
\]
Its units on a common chart are the units of $w$ pulled
along $\tau_x$, divided by those pulled along
$\operatorname{pr}_{A'}$.  Thus (PBC.49) is a comparison
of the literal $\operatorname{translDiff}$ bundles.

Now suppose the antecedent in $\operatorname{KTriv}(A',L')$
holds.  For every $p\in T$ it gives an open $V\ni p$ and
an actual generating tuple of the left side of (PBC.49)
over $\operatorname{pr}_T^{-1}V$.
The isomorphism $\alpha$ commutes with $\operatorname{pr}_T$,
so it restricts to an isomorphism over this same $V$.
Transport the tuple by (PBC.49) and by this restricted
isomorphism.  It gives an actual generator of
$\operatorname{translDiff}(A,L,x_0)$ over
$\operatorname{pr}_T^{-1}V$ in $Y$.
Consequently $x_0$ satisfies the antecedent of the
original $\operatorname{KTriv}(A,L)$ for the same ring
$R$ and the composite map $f\circ t$.
That condition forces $x_0=e_A\circ f\circ t$.

PBC22 identifies this latter point as the image of the
identity point of $A'(T)$, and injectivity there gives
$x=e_{A'}\circ t$.  This is the exact original conclusion.
The argument used an isomorphism in (PBC.47), not any
unproved descent of sections along a nonflat map.
\end{proof}
% END PBC24

% BEGIN PBC25
\paragraph{PBC25: full polarized base change with all original fields.}
\label{claim:PBC25}
For every pair of schemes $S,S'$, every $f:S'\to S$, and
every original polarized record $P=\operatorname{PolAV}(S)$,
there are a polarized record $P'$ over $S'$ and a Cartesian
map $\phi:P'.A\to P.A$ over $f$ such that
\[
 \begin{gathered}
 P'.g=P.g,\\
 \operatorname{toBase}(x y)
       =\operatorname{toBase}(x)\operatorname{toBase}(y)
       \quad\text{for every }T/S'\text{ and every }x,y\in P'.A(T),\\
 \text{there is an actual generated }
       \operatorname{LB.Iso}(P.L.\operatorname{pull}(\phi),P'.L).
 \end{gathered}                                           \tag{PBC.50}
\]
In fact $P'.A$ is the specified categorical pullback
$P.A\times_S S'$, $P'.L$ is literally $P.L.\operatorname{pull}(\phi)$,
and $P'.g$ is defined to be $P.g$.
This proves the full original $\operatorname{PolBaseChangeSt}$.

\begin{proof}
Use the actual pullback from PBC22 with its pulled group
structure.  It is abelian and has the same smooth relative
dimension $g=P.g$.  Set $L'=\phi^*P.L$ with the original
pullback charts.  Pulling the transition units preserves
the cocycle and the covering, so $L'$ is an actual line bundle.
PBC23 supplies its rigidification and local symmetry from
$P.\operatorname{rigid}$ and $P.\operatorname{sym}$.
PBC24 supplies its $\operatorname{KTriv}$ from
$P.\operatorname{ktriv}$, using the identity comparison
on $L'$ constructed in PBC19.

It remains to check the relative-ample field on this same
bundle.  The pullback is proper by PBC22, hence its
structural morphism is quasi-compact and quasi-separated.
For every $s'\in S'$, choose an affine open $V\subseteq S$
containing $f(s')$.  PBC21 gives relative ampleness of the
literal restriction of $L'$ over $f^{-1}V$.
PBC18 then proves relative ampleness on all of $S'$.

These constructions supply every field of the original
$\operatorname{PolAV}$ record: group, abelian property,
$g$, smooth relative dimension, actual bundle, line-bundle
compatibility, rigidification, symmetry, relative ampleness,
and $\operatorname{KTriv}$.  Define $P'$ using them.
The Cartesian square is the original pullback square.
The dimension equality in (PBC.50) is equality by the
definition of $P'.g$.  The universal multiplication assertion
is exactly PBC22, for every test $T$, without an affine or
field restriction on that test.

Finally, since $P'.L=P.L.\operatorname{pull}(\phi)$,
use the actual identity division-bundle generator from
PBC19.  Its coordinate on the pair chart $(i,j)$ is the
pulled bundle's transition $g^{L'}_{ij}$; it satisfies
the generator compatibility by the cocycle identity.
Thus the last output is the source's generated
$\operatorname{LB.Iso}$, not only an isomorphism class
of line bundles.  This completes (PBC.50) for arbitrary
$S,S'$ and $f$.
\end{proof}
% END PBC25

% BEGIN PBC26
\paragraph{PBC26: the complete all-affine-test very-ample producer.}
\label{claim:PBC26}
For every original $P/S$, every commutative ring $R$, every
$t:\operatorname{Spec}R\to S$, and every specified Cartesian
model
\[
 \begin{array}{ccc}
 X&\xrightarrow{\pi}&P.A\\
 f_X\downarrow&&\downarrow\\
 \operatorname{Spec}R&\xrightarrow{t}&S,
 \end{array}
\]
there are $N\ge1$, $m\ge0$ and actual sections
$s_0,\ldots,s_m$ of
$\operatorname{pw}(P.L.\operatorname{pull}(\pi),N)$
such that this power is an actual line bundle and
$\operatorname{VeryAmpleBy}(f_X,\operatorname{pw}(\pi^*P.L,N),s)$.
This is the original \texttt{polAV\_veryAmpleBy} with its
$\operatorname{PolBaseChangeSt}$ premise now supplied.

\begin{proof}
PBC25 provides the polarized pullback over
$\operatorname{Spec}R$.  The prescribed $X$ is uniquely
isomorphic to its underlying pullback, commuting with both
projections.  Apply its relative ampleness to the whole affine
base.  It gives one positive power, the finite tuple of
sections, affine nonzero loci, actual ratio identities and
the surjective polynomial maps over
$\mathcal O(\operatorname{Spec}R)$.

The already proved GSL15 transports precisely these data
through the unique Cartesian isomorphism and the generated
bundle comparison to $X$.  Its argument uses GSL14 for
the actual sections, PBC19's unit cancellation for the
ratio functions, and the canonical coefficient-ring
isomorphism $R\simeq\mathcal O(\operatorname{Spec}R)$ for
the polynomial maps.  It also changes the repeated-factor
power to the literal $\operatorname{pw}$ chart record.
The resulting data are exactly the asserted
$\operatorname{VeryAmpleBy}$, including the line-bundle
compatibility of that power.

The quantifiers allow $N$ and the tuple to depend on this
specified affine test and Cartesian model.  They do not
assert one uniform ample power over all tests.
The ring $R$ is arbitrary, with no reducedness, locality,
Noetherianity or flatness condition.
\end{proof}
% END PBC26

% BEGIN PBC27
\paragraph{PBC27: actual section finiteness and the return to the original power-section module.}
\label{claim:PBC27}
Let $K$ be any field, $f:X\to\operatorname{Spec}K$, $M$
an original line-bundle chart datum, and
$Q=\operatorname{Perf}(f,M)$.
If the degree-zero matrix size of $Q$ is $a$, every
$a+1$ actual sections $s_0,\ldots,s_a$ satisfy a nonzero
$K$-linear relation in their coordinates on every original
$M$-chart.
Consequently, for every algebraically closed $K$, every
original $P/\operatorname{Spec}K$ and every actual line
bundle $M$ on $P.A$, there is one $n$ such that every
$n$-tuple of actual $M$-sections has a nonzero relation
over $\mathcal O(\operatorname{Spec}K)$, using the original
structural maps on every chart.  Moreover
\[
 \operatorname{SecM}(P.A,P.L,j,\operatorname{baseOver}(\operatorname{Spec}K))
 \quad\hbox{is finite over }\mathcal O(\operatorname{Spec}K)
       \quad\hbox{for every }j\in\mathbb Z.
                                                        \tag{PBC.51}
\]
More generally, over any field $K$, conclusion (PBC.51)
follows whenever the displayed finite-relation property is
supplied for every actual line bundle on $P.A$; algebraic
closedness is not a premise of this last implication.

\begin{proof}
Use the identity base change in $Q.\operatorname{sec}$.
For each $s_k$, first express it on the literal identity-pullback
charts of $M$.  Its preimage under the provided section
bijection is a vector $v_k\in\ker D_0\subseteq K^a$.
The actual coordinate identity in $\operatorname{Perf}$ says
that on its fixed cover $W_\alpha\subseteq M.U_{c(\alpha)}$,
\[
 s_{k,c(\alpha)}|_{W_\alpha}
    =\sum_{\ell=1}^{a} f^*(v_{k\ell})u_{\ell\alpha}.
                                                        \tag{PBC.52}
\]
This identity, rather than any unstated linearity of the
chosen bijection, is all that is needed.
The $a+1$ vectors in $K^a$ have a relation
$\sum_k d_kv_k=0$ with $(d_k)\ne0$.
For $a=0$, the sole vector is zero and $d_0=1$ works.
Multiply (PBC.52) by $f^*d_k$, sum in $k$, and interchange
the finite sums.  Every coefficient of every
$u_{\ell\alpha}$ is zero.  Thus
$\sum_k f^*(d_k)s_{k,c(\alpha)}=0$ on $W_\alpha$.

For any original chart $M.U_i$, restrict to
$M.U_i\cap W_\alpha$ and use the transition unit
$g_{i,c(\alpha)}$ to transfer this same relation to the
$i$th coordinates.  The opens $W_\alpha$ cover $X$.
Equality of sections on a cover proves
$\sum_k f^*(d_k)s_{k,i}=0$ on all of $M.U_i$.
This is the full original $\operatorname{Perf.exists\_relation}$.

For the asserted polarized variety, use PBC26 on the
identity Cartesian model of $P.A$.  It supplies the
original $\operatorname{IsProjBy}$ presentation.
The structure map is smooth by $P.\operatorname{isAb}$.
ASL21 therefore provides an actual
$Q=\operatorname{Perf}(P.A\to\operatorname{Spec}K,M)$.
Choose $n=Q.a+1$ before the tuple of sections.
Apply the first part and send its coefficients through the
canonical ring isomorphism
$K\simeq\mathcal O(\operatorname{Spec}K)$.
It is injective, so the coefficient tuple remains nonzero.
Its composite with the structural map to each chart is
exactly the scalar map in (PBC.52).
This proves the complete original
\texttt{sections\_finite\_of\_perf} conclusion.

To obtain (PBC.51), put $T=\operatorname{baseOver}(\operatorname{Spec}K)$.
The second projection
$\beta:T\times_{\operatorname{Spec}K}P.A\to P.A$
is an isomorphism with inverse the pair of the structural
map and the identity of $P.A$; the two projection identities
prove this directly.
On $P.A$ take the actual integer-power bundle
$M_j=\operatorname{LinComb}(P.L,\operatorname{id},j)$.
Apply the just proved finite relation statement to $M_j$.
A tuple in the source's $\operatorname{SecM}$ is a tuple
of actual sections on the identity fiber product.
Pull it back by $\beta^{-1}$, identifying its finite-singleton
power charts with those of $M_j$ by PBC19.
The resulting nonzero relation pulls back to the same
relation on the fiber product, because the scalar maps
commute with both projections.  Thus every tuple of the
chosen length in $\operatorname{SecM}$ is linearly dependent.

Finally, over a field, if every $n$-tuple is dependent then
the whole vector space is finitely generated: starting with
the empty independent family, repeatedly adjoin a vector
outside its span; this either terminates with a spanning
family or would produce an independent $n$-tuple.
Here the coefficient ring is a field by
$K\simeq\mathcal O(\operatorname{Spec}K)$.
This proves finiteness with the original coefficient action.
The identity-product argument works over any field whenever
the displayed original finite-relation hypothesis is supplied,
which is also the precise scope of the source's
\texttt{finite\_secM}.
\end{proof}
% END PBC27

\subsection{Return to the section-law construction and exact remaining scope}

PBC25 supplies the original input denoted $\mathcal P$ in GSL15--GSL17
and \texttt{GSL\_POL\_BASE\_CHANGE} in the existing contracts.
The construction retains a fixed $P$ and pulls its actual group,
line bundle and specified generators.  It therefore can be inserted
in the previous conditional arguments without changing their sections
or their coefficient maps.
PBC26 supplies their all-test projective presentation.
ASL21's original smooth-projective perfect-complex theorem can now
be applied with that presentation on the specified affine pullback.
The completed ASL and SHV proofs are not repeated.

The two Serre conclusions of SHV17 and SHV21 remain available at their
original Noetherian-base scopes.  They do not impose Noetherianity on
PBC18 or PBC25, which were proved independently over arbitrary bases.
Conversely, supplying polarized base change does not prove the
remaining geometric inputs to GSL17: the multiplication-pullback
generator, the finite/flat/surjective and away-characteristic etale
multiplication geometry, and invariant-function descent on every
preimage of an open still have to be supplied.

PBC27 supplies actual section finiteness, including every integer
power in the original $\operatorname{SecM}$ model.
The geometric rank-one theorem, torsion-indexed section decomposition,
two-sided section growth, and the square/translation generators are also
separate original producers.  Until their exact remaining hypotheses
are discharged, full global
$\operatorname{SecBaseChangeSt}$,
$\operatorname{SecBaseChangeAnySt}$ and
$\operatorname{SecTwistSt}$ are not asserted by this chapter.
No change is made to the retained moduli or finite-level dictionary
and class-sum obligations, and no whole selected obligation is closed.

\paragraph{Source map.}
The following references identify the mathematical producer bodies;
the accompanying provenance binds their exact bytes and line ranges.
PBC1--PBC9 reconstruct
\texttt{Exists\allowbreak One\allowbreak Add\allowbreak Kill.lean}: the regular-function calculus,
\texttt{kill}, \texttt{frac}, \texttt{local\_global},
\texttt{exists\_one\_add\_kill}, \texttt{Sc.extend}, and
\texttt{extra\_section}.
PBC10--PBC18 reconstruct
\texttt{Rel\allowbreak Ample\allowbreak Of\allowbreak Locally\allowbreak On\allowbreak Base.lean}: the full
\texttt{glue\_pack}, open-immersion transport,
\texttt{gen\_iff}, the actual-section dictionary,
\texttt{exists\_sv}, \texttt{cst\_mem},
\texttt{pack\_of\_relAmple}, and
\texttt{relAmple\_of\_locally}.

PBC19--PBC24 reconstruct the relevant complete chart and pullback
arguments in \texttt{AmplePt.lean}, \texttt{SabPull.lean},
\texttt{ThetaCore.lean}, \texttt{IsMonHomOfMul.lean},
\texttt{ConstPullback.lean}, and
\texttt{Pol\allowbreak Base\allowbreak Change\allowbreak Of\allowbreak Rel\allowbreak Ample\allowbreak Local.lean}.
The last of these supplies the constructor returned by PBC25.
The literal polarization and all-test output contracts are in
\texttt{NonzeroLocus.lean} and
\texttt{GlobalSectionsRankAlgClosed.lean}.
PBC26 returns to
\texttt{VeryAmpleByOfRelAmple.lean}, using the already reviewed
GSL14--GSL15 proof blocks at their exact scopes.
PBC27 proves the original producers in
\texttt{AbelianVarietyLineBundleSectionsFinite.lean} and
the \texttt{finite\_secM} argument of
\texttt{FiniteOfNotLinearIndependent.lean}, using ASL21.
